% ==============================================================================
% VHP-02 — Scientific Paper Template (arXiv Style)
% Venture Hackathon Package (VHP-001) / Mindrock Research
% Canonical Architecture: Conditional Theorems, Countermodels & Axiom Ledger
% ==============================================================================
\documentclass[11pt,a4paper]{article}

% --- Core Typography & Geometry ---
\usepackage[utf8]{inputenc}
\usepackage[T1]{fontenc}
\usepackage{lmodern}
\usepackage[margin=1in]{geometry}
\usepackage{microtype}
\usepackage{amsmath,amssymb,amsthm,amsfonts}
\usepackage{booktabs}
\usepackage{tabularx}
\usepackage{cite}
\usepackage{xcolor}
\usepackage{tcolorbox}
\usepackage{listings}
\usepackage{hyperref}
\usepackage[capitalise,noabbrev]{cleveref}

% --- Hyperlink Setup ---
\hypersetup{
    colorlinks=true,
    linkcolor=blue!70!black,
    citecolor=green!50!black,
    urlcolor=blue!80!black
}

% --- Theorem Environments ---
\theoremstyle{plain}
\newtheorem{theorem}{Theorem}[section]
\newtheorem{conditionaltheorem}[theorem]{Conditional Theorem}
\newtheorem{proposition}[theorem]{Proposition}
\newtheorem{lemma}[theorem]{Lemma}
\newtheorem{corollary}[theorem]{Corollary}

\theoremstyle{definition}
\newtheorem{definition}[theorem]{Definition}
\newtheorem{axiom}[theorem]{Axiom}
\newtheorem{assumption}[theorem]{Assumption}
\newtheorem{countermodel}[theorem]{Countermodel}

\theoremstyle{remark}
\newtheorem{remark}[theorem]{Remark}
\newtheorem{example}[theorem]{Example}

% --- Boxed Claim Environments ---
\tcbuselibrary{skins,breakable}
\newtcolorbox{claimbox}[2][]{
    enhanced,
    colback=gray!5,
    colframe=blue!40!black,
    fonttitle=\bfseries,
    title={#2},
    arc=1mm,
    boxrule=0.5mm,
    breakable,
    #1
}

\newtcolorbox{counterbox}[2][]{
    enhanced,
    colback=red!5,
    colframe=red!60!black,
    fonttitle=\bfseries,
    title={#2},
    arc=1mm,
    boxrule=0.5mm,
    breakable,
    #1
}

% --- Notation Macros ---
\newcommand{\lean}[1]{\texttt{#1}}
\newcommand{\artifact}[1]{\texttt{#1}}
\newcommand{\sha}[1]{\texttt{#1}}
\newcommand{\cmd}[1]{\texttt{#1}}

% --- Title & Metadata ---
\title{\textbf{Title of the Paper: Subtitle Indicating Formal Scope}\\[0.5ex]
\large \textsc{VHP-02 arXiv-Style Formal Paper Template}}

\author{
    \textbf{Author One}\thanks{Equal contribution. Mindrock Autonomous Network. Email: \texttt{author1@mindrock.capital}}
    \and
    \textbf{Author Two}\thanks{Equal contribution. Research Lab. Email: \texttt{author2@mindrock.capital}}
    \and
    \textbf{Pavel Cherkashin}\thanks{Principal Sponsor and Architecture Authority. Mindrock Capital. Email: \texttt{pavel@mindrock.capital}}
}

\date{\today}

\begin{document}

\maketitle

\begin{abstract}
\noindent
State the exact core result in 150--250 words. The abstract must clearly separate:
(1) what is formally proved and kernel-checked;
(2) what is conditional upon explicit hypotheses;
(3) what classical properties are affirmatively falsified with concrete countermodels; and
(4) where the exact verification artifacts and reproducible replay commands are located.
Every scientific claim asserted in this paper is indexed in the machine-readable traceability matrix.
\end{abstract}

\paragraph{Keywords:} Formal Verification, Conditional Theorems, Countermodels, Axiom Ledger, Lean 4, Reproducibility.
\paragraph{MSC (2020):} Primary 08A05; Secondary 68V15, 03B35.
\paragraph{ACM CCS:} Theory of computation $\to$ Logic and verification; Formal languages.

\vspace{1em}
\hrule
\vspace{1em}

% ==============================================================================
% 1. INTRODUCTION
% ==============================================================================
\section{Introduction}
\label{sec:intro}

Provide the motivation, the historical tension between informal proofs and formal rigor, and the necessity of separating unconditional deductions from conditional statements.

\subsection{Contributions}
\begin{enumerate}
    \item \textbf{Axiom Ledger (\cref{sec:axiom-ledger}):} An explicit, bounded catalog of foundational axioms and domain constraints.
    \item \textbf{Conditional Theorems (\cref{sec:conditional-theorems}):} Formal results whose truth is explicitly parameterized by listed premises.
    \item \textbf{Countermodels \& Falsification (\cref{sec:countermodels}):} Explicit counterexamples demonstrating that naive classical extensions fail.
    \item \textbf{Traceability Matrix (\cref{sec:traceability}):} A bidirectional mapping from every textual claim to its formal artifact, hash receipt, and replay command.
\end{enumerate}

% ==============================================================================
% 2. FORMAL DEFINITIONS & AXIOM LEDGER
% ==============================================================================
\section{Formal Definitions and Axiom Ledger}
\label{sec:axiom-ledger}

State all algebraic structures and primitives. The Axiom Ledger categorizes every assumption into:
\begin{itemize}
    \item \textbf{Foundational Axioms:} Trusted axioms of the underlying logical framework (e.g., Lean 4 kernel: propositional extensionality, quotient soundness, classical choice).
    \item \textbf{Domain Axioms:} Defining postulates for the mathematical theory under study.
    \item \textbf{Conjectural Postulates:} Explicit conditional hypotheses that remain unproved.
\end{itemize}

\begin{table}[htbp]
\centering
\small
\caption{Axiom Ledger Summary}
\label{tab:axiom-ledger}
\begin{tabularx}{\textwidth}{llllX}
\toprule
\textbf{Axiom ID} & \textbf{Category} & \textbf{Name} & \textbf{Kernel Status} & \textbf{Formal Statement} \\
\midrule
AXM-01 & Foundational & Equality Reflexivity & Lean Kernel & $\forall x, x = x$ \\
AXM-02 & Domain & Involution of Reciprocal & Proved / Admitted & $\forall x, /(/x) = x$ \\
AXM-03 & Domain & Regularized Distributivity & Proved & $(x+y)z + 0z = xz + yz$ \\
\bottomrule
\end{tabularx}
\end{table}

\begin{definition}[Core Structure]
\label{def:structure}
Let $(W, +, \cdot, /, 0, 1)$ be a structure equipped with binary operations addition $(+)$ and multiplication $(\cdot)$, a unary reciprocal $(/)$, and distinguished constants $0$ and $1$.
\end{definition}

% ==============================================================================
% 3. CONDITIONAL THEOREMS
% ==============================================================================
\section{Conditional Theorems}
\label{sec:conditional-theorems}

A theorem is strictly defined as \textit{conditional} when its conclusion holds only in the presence of specified hypotheses. Each theorem statement in this section explicitly names its required subset of the Axiom Ledger.

\begin{claimbox}{Conditional Theorem \ref{thm:sample-conditional}: Title of Theorem}
\begin{conditionaltheorem}[Dependent on Axioms AXM-01, AXM-02]
\label{thm:sample-conditional}
Let $W$ satisfy Axioms AXM-01 and AXM-02. Then for all non-singular elements $x, y \in W$, property $\mathcal{P}(x, y)$ holds.
\end{conditionaltheorem}
\end{claimbox}

\begin{proof}
Proof steps referencing declarations in the formal library. Each deduction step is verified by the proof assistant kernel.
\end{proof}

\begin{proposition}[Invariant Regularization]
\label{prop:regularization}
If hypothesis $\mathcal{H}_1$ is satisfied, then operator $\mathcal{T}$ regularizes singularities without collapsing the non-trivial equivalence classes.
\end{proposition}

% ==============================================================================
% 4. COUNTERMODELS & BOUNDARY FALSIFICATION
% ==============================================================================
\section{Countermodels and Boundary Falsification}
\label{sec:countermodels}

To prevent false claims of generality, this section catalogues explicit boundary witnesses where classical laws fail. Discovering and recording countermodels is a primary scientific contribution.

\begin{counterbox}{Countermodel \ref{cm:sample-counter}: Refutation of Classical Property}
\begin{countermodel}[Falsification of Unrestricted Property $\mathcal{Q}$]
\label{cm:sample-counter}
The property $\forall x \in W, x - x = 0$ is \textbf{false} in the general structure.
\begin{itemize}
    \item \textbf{Witness:} Let $x = \infty$ (where $\infty = /0$).
    \item \textbf{Evaluation:} $\infty - \infty = \infty + (-\infty) = \bot \neq 0$.
    \item \textbf{Conclusion:} Classical additive annihilation fails at singularities.
\end{itemize}
\end{countermodel}
\end{counterbox}

\begin{table}[htbp]
\centering
\small
\caption{Countermodel Falsification Ledger}
\label{tab:countermodels}
\begin{tabularx}{\textwidth}{llXX}
\toprule
\textbf{Countermodel ID} & \textbf{Target Classical Law} & \textbf{Status} & \textbf{Explicit Witness \& Refutation} \\
\midrule
CM-01 & Additive Inverse ($x-x=0$) & FALSIFIED & $x=\infty \implies \infty-\infty = \bot \neq 0$ \\
CM-02 & Zero Annihilation ($0 \cdot x = 0$) & FALSIFIED & $x=\infty \implies 0 \cdot \infty = \bot \neq 0$ \\
CM-03 & Classical Distributivity & FALSIFIED & $x=1, y=1, z=\infty \implies (1+1)\infty \neq 1\infty + 1\infty$ \\
\bottomrule
\end{tabularx}
\end{table}

% ==============================================================================
% 5. VERIFICATION ARTIFACT TRACEABILITY MATRIX
% ==============================================================================
\section{Verification Artifact Traceability Matrix}
\label{sec:traceability}

In adherence to the Mindrock Handoff Gate Protocol (MIN-236 / AGENTS.md), every claim asserted in this text must be mechanically traceable to an immutable code artifact.

\begin{table}[htbp]
\centering
\scriptsize
\caption{Claims Traceability Matrix}
\label{tab:traceability}
\begin{tabularx}{\textwidth}{llllX}
\toprule
\textbf{Claim / Theorem} & \textbf{Artifact Path} & \textbf{Declaration / Symbol} & \textbf{Audit Status} & \textbf{Verification Command} \\
\midrule
\cref{thm:sample-conditional} & \artifact{src/Core/Wheel.lean} & \lean{Wheel.conditional\_thm} & Kernel Checked & \cmd{lake build && lean --run audit.lean} \\
\cref{prop:regularization} & \artifact{src/Algebra/Reg.lean} & \lean{Wheel.regularize\_invol} & Zero-Sorry Gate & \cmd{python3 -m unittest tests/test\_reg.py} \\
\cref{cm:sample-counter} & \artifact{tests/Countermodels.rs} & \cmd{test\_countermodel\_01} & Replay Passed & \cmd{cargo test --test countermodels} \\
\bottomrule
\end{tabularx}
\end{table}

% ==============================================================================
% 6. COMPUTATIONAL VERIFICATION & KERNEL AUDIT
% ==============================================================================
\section{Computational Verification and Kernel Audit}
\label{sec:verification}

Describe the formal proof engine environment:
\begin{itemize}
    \item \textbf{Proof Engine:} Lean 4 pinned toolchain (e.g. \texttt{leanprover/lean4:v4.14.0}) or Rust formal check suite.
    \item \textbf{Axiom Audit:} Execution of \texttt{\#print axioms} confirming zero unreviewed axioms (\texttt{sorry}, \texttt{admit}, or unsafe primitives).
    \item \textbf{Deterministic Build:} Pinned lockfile ensuring bitwise identical execution across independent runs.
\end{itemize}

% ==============================================================================
% 7. LIMITATIONS AND OPEN QUESTIONS
% ==============================================================================
\section{Limitations and Open Questions}
\label{sec:limitations}

State with transparency what has \textbf{not} been established:
\begin{enumerate}
    \item Boundary limitations of the algebraic model.
    \item Unsettled conjectures not covered by the current kernel checks.
    \item Computational complexity bounds of arbitrary equality checks.
\end{enumerate}

% ==============================================================================
% 8. REPRODUCIBILITY & INDEPENDENT REPLAY
% ==============================================================================
\section{Reproducibility and Independent Replay}
\label{sec:reproducibility}

An independent verifier can replicate all proofs, tests, and countermodels with the following zero-trust checkout command:

\begin{lstlisting}[language=bash,basicstyle=\ttfamily\small,breaklines=true,frame=single]
# Clone from canonical remote repository
git clone --branch <branch-name> <canonical-repo-url> verify-box
cd verify-box
git checkout <commit-sha>

# Verify integrity and replay verification suite
python3 tools/verify-vhp-paper verify
\end{lstlisting}

\bibliographystyle{plain}
\begin{thebibliography}{99}
\bibitem{carlstrom2004}
J.~Carlstr{\"o}m.
\newblock Wheels -- on division by zero.
\newblock {\em Mathematical Structures in Computer Science}, 14(3):403--436, 2004.

\bibitem{demoura2021}
L.~de~Moura and S.~Ullrich.
\newblock The {L}ean 4 theorem prover and programming language.
\newblock In {\em Automated Deduction -- CADE 28}, pages 625--635. Springer, 2021.
\end{thebibliography}

\end{document}
